% Drafted 2026-08-17 against the elsarticle version of the manuscript.
% Inputted by paper_journal/main.tex immediately after the feature ablation,
% because it interprets that result rather than standing alone.
% Numbers trace to experiments/results/final_eval/shap_*.csv and to the
% RESEARCH_LOG.md entry dated 2026-08-13.

\subsection{Attribution: What the Router Actually Uses}
\label{sec:shap}

The ablation in Section~\ref{sec:ablation} establishes that patient features
carry no \emph{performance}: removing them costs nothing. It does not establish
that the router ignores them. Those are different claims, and on this router
they come apart.

We compute Shapley additive explanations~\cite{lundberg2017} over the deployed
classifier's real 389-dimensional input, using the exact fitted embedding and
scaler from the serialised pipeline rather than a reconstruction. Two
correctness gates precede any attribution: on the routing-validation split the
recomputed class probabilities match the stored predictions to $3\times10^{-8}$,
and on the held-out evaluation set the recomputed routing decisions reproduce
all 300 deployed mode selections exactly. Attributions satisfy the additivity
identity to within $8\times10^{-7}$ of the raw model margin. Without those
checks an explanation may describe a feature vector the deployed router never
saw.

\subsubsection*{The block view corroborates the ablation}

Aggregating mean $|\phi|$ within each block, the 384 question-embedding
dimensions account for \textbf{91.2\%} of total attribution mass and the five
structural patient features for \textbf{8.8\%} (Fig.~\ref{fig:shap}a). Per
class the patient share is 9.0\% for text-only routing, 9.0\% for text with
EHR, and 8.4\% for the full hybrid mode --- the knowledge-graph decision is the
\emph{least} patient-informed of the three, which is consistent with our
sparsity finding that graph retrieval is keyed on the diagnosis list rather
than on patient state.

\subsubsection*{The per-feature view does not}

That block comparison sums 384 features against five, which flatters the
larger block by construction. Per feature the ordering reverses: the average
patient feature receives \textbf{7.4 times} the attribution of the average
embedding dimension (Fig.~\ref{fig:shap}b). Ranked across all 389 inputs,
\texttt{n\_labs} is the \textbf{fifth} most attributed feature and is non-zero
on all 300 held-out decisions; \texttt{n\_diag} ranks 14th and
\texttt{sparsity\_score} 17th. Only \texttt{n\_meds} is genuinely inert, with
attribution exactly zero on every decision --- a dead input rather than a weak
one.

So the router demonstrably consults patient state on every decision it makes.
The ablation shows that consulting it does not improve routing. Both are true,
and the defensible characterisation is a \textbf{question-driven policy that
makes performance-neutral use of patient state} --- not a policy that ignores
it. We stress the distinction because the weaker reading is easy to reach from
either measurement alone, and it would overstate a negative result that is
already sufficiently strong: patient features are redundant here, not unused.

The most plausible reconciliation is redundancy rather than irrelevance. On
laboratory questions \texttt{n\_labs} separates the two viable modes perfectly
--- admissions with no recorded laboratory values receive text-only routing ---
but the question text already signals the same distinction, so removing the
feature costs no accuracy while its attribution remains high. A single
held-out decision illustrates the pattern: for a laboratory question on an
admission with no recorded results, \texttt{n\_labs} is the highest-attributed
of all 389 inputs, contributing $+0.81$ toward the text-only decision.

\subsubsection*{The attributions are causal for this model}

Attribution magnitude alone does not establish that a feature drives the
decision. We test this by deletion: for each decision we ablate the $k$
features SHAP credits most toward the predicted class, replacing them with
their training-set means, and measure the change in that class's probability
against a random-$k$ control at the same budget (Fig.~\ref{fig:shap}c).

Ablating SHAP-ranked features reduces the predicted-class probability in
\textbf{89.9\%} of cases overall and in \textbf{100\%} once $k \geq 10$, and
flips the routing decision outright in 25.5\% of cases. The random control
moves the model barely at all: area over the deletion curve is \textbf{25.8
times} larger for SHAP-ranked features than for randomly chosen ones, with a
decision flip rate of 1.1\% against 25.5\%. Both curves plateau by $k \approx
20$, indicating that routing rests on a small subset of the 389 inputs rather
than on diffuse mass across them.

We note the standard caveat: replacing a dense embedding dimension with its
training mean is an off-manifold perturbation, so this measures the model's
local sensitivity to the features SHAP credits, not the effect of a realisable
intervention.

\begin{figure}[t]
\centering
\includegraphics[width=\linewidth]{shap_paper_figure_tierB.png}
\caption{SHAP attribution over the deployed router's 389 inputs on the 300
held-out routing decisions. (a) By block, question embeddings dominate,
corroborating the ablation. (b) Per feature the ordering reverses: four of the
five structural patient features sit above the mean embedding dimension
(dashed), with \texttt{n\_labs} ranked fifth of 389; \texttt{n\_meds} has
attribution exactly zero and is plotted at the origin. (c) Deletion test:
ablating SHAP-ranked features collapses the predicted-mode probability while a
random-$k$ control at the same budget does not.}
\label{fig:shap}
\end{figure}
